\chapter{资产、负债、净值：全世界通用的三个词}
\chaptermeta{05}

\begin{ledetext}
上一章说通胀在债权人和债务人之间搬钱。想看清怎么搬，得先看清账是怎么记的。这一章给的不是知识，是一件工具。
\end{ledetext}

\section{王雪梅的三个数字}

王雪梅 31 岁，合肥人。2025 年秋天教培机构裁员，她那个课程顾问的岗位整组没了。丈夫周立在一家汽车电子公司做测试工程师。

失业第三个月的一个晚上，她把银行 App、公积金页面和两张还款计划表一起摊在餐桌上，第一次把家里的账算了一遍。

算法只用到三个词。

\term{资产}{assets}\index{093@资产}，你拥有的、以后能给你带来好处的东西：房子、存款、基金、一台还能干活的机器，还有别人欠你的钱。\term{负债}{liabilities}\index{017@负债}，你欠别人的、以后必须交出去的东西：房贷、车贷、信用卡账单、公司欠供应商的货款。\term{净值}{net worth}\index{030@净值}，两者之差。

\begin{pullquote}{}
资产 = 负债 + 净值
\end{pullquote}

它永远成立，是因为净值不是量出来的，是\textbf{被定义成残值}的：把你有的数清楚，把你欠的减掉，剩多少算多少。没有作弊的空间，也没有解释的余地。

\begin{bsheet}
  \bsside{资产 ASSETS}{%
    \bsrow{滨湖的房子（同户型挂牌价）}{232 万}
    \bsrow{存款}{18.6 万}
    \bsrow{两人的公积金账户}{9.4 万}
    \bsrow{车（二手估价）}{7.3 万}
    \bsrow{基金}{4.1 万}
    \bsrowtotal{合计}{271.4 万}
  }
  \bsside{负债 + 净值}{%
    \bsrow{房贷余额}{176.4 万}
    \bsrow{车贷余额}{2.8 万}
    \bsrow{信用卡待还}{1.9 万}
    \bsrow{净值}{90.3 万}
    \bsrowtotal{合计}{271.4 万}
  }
\end{bsheet}

这套房 2021 年买进来的时候是 268 万。

光是挂牌价从 268 万滑到 232 万这一件事，就从这张表的净值里拿走了 36 万。而房贷余额那一栏，一分钱都没跟着少。

这个不对称本章结尾还要再算一次。

左右两栏必须严丝合缝地相等。不等，就是有东西没数清楚。

公司的表长得一模一样，只是"净值"那一栏换了个名字，叫\term{所有者权益}{equity}\index{057@所有者权益}。同一个位置，同一套算法：资产数清楚，欠人的减掉，剩下的归股东。第 6 章讲"股"这张合同的时候，讲的就是这一栏。

\section{五百年前的记账法}

这套东西叫\term{复式记账}{double-entry bookkeeping}\index{019@复式记账}。规则只有一句：任何一笔交易都会同时改动至少两个科目，改完之后等式照样成立。

1494 年，方济各会修士卢卡·帕乔利在威尼斯出版了一本数学书，顺手把当地商人已经在用的这套办法写成了系统。五百多年过去，从王雪梅手机上那行流水到美联储的资产负债表，用的还是同一套语法。

它的本事在于，它逼你回答"另一边发生了什么"。

\begin{egbox}{举个栗子}
王雪梅想拿 11 万提前还房贷，理由是"这样心里踏实"。查账：

资产端存款 18.6 万变成 \textbf{7.6 万}，负债端房贷 176.4 万变成 \textbf{165.4 万}。

新合计：资产 260.4 万 = 负债 170.1 万 + 净值 \textbf{90.3 万}。

\textbf{净值一分没动。}

提前还贷让表的两边同时缩小，它改的是往后的利息支出和风险敞口，不是财富。

对照一下：周立年底要是拿到 4.2 万奖金并存下来，资产加 4.2 万，负债不动，净值从 90.3 万变 \textbf{94.5 万}。这才叫变富。

"心里踏实"是个正当理由。只是它买的不是财富，是睡眠。

\end{egbox}

\section{全书最重要的一句话}

\begin{pullquote}{}
一个人的金融资产，一定是另一个人的负债。
\end{pullquote}

王雪梅手机里那 18.6 万余额，她以为是"她的钱"。在银行的账本上，它躺在\textbf{负债}那一栏，是银行欠她的。银行拿什么撑住这笔负债？资产栏里的贷款，也就是别人欠银行的。

\begin{chainflow}
\chainnode{王雪梅的 18.6 万存款}\chainarrow \chainnode{是银行的负债}\chainarrow \chainnode{银行的资产是贷款}\chainarrow \chainnode{那是别人的负债}\chainarrow \chainnode{别人的资产是房子和厂房}\chainarrow \chainnode{实物资产，链条到头}
\end{chainflow}

链条可以拉得很长：她的存款、银行的债券、企业的应付账款、政府的国债。但它必须终止在两处之一，要么是实物资产，要么是股权。股权是对剩余收益的索取权，不是一张固定金额的欠条。

换句话说，金融本身不生产任何东西。它只把实物世界的收益重新切开，分给不同的人和不同的时间。第 6 章会把这句话变成整本书的支点。

\begin{databox}{2026 数据}
国际金融协会（IIF）的统计：截至 \textbf{2026 年 3 月底}，全球债务总额约 \textbf{353 万亿美元}，创历史新高，相当于全球 GDP 的约 \textbf{305\%}。同期，美国国债规模在 2026 年 8 月\textbf{突破 40 万亿美元}。

标题很吓人。把纸翻过来看背面：这 353 万亿同时是另一些人的 \textbf{353 万亿资产}。养老金账户里的国债，保险公司的准备金，银行资产栏里的贷款，王雪梅那 18.6 万。

所以债务的\emph{总量}本身说明不了危险。危险在结构里：谁欠的、欠谁、什么币种、什么期限、拿什么现金流还。第 14 章和第 17 章拆这个。

\end{databox}

\section{谁更有钱}

\begin{egbox}{举个栗子}
\textbf{李振宇}，34 岁，周立的同事，做项目经理，月薪 3 万 2，到手两万四出头。

资产：高新区的房子现在挂 415 万、车 22 万、存款 3.2 万，合计 \textbf{440.2 万}。负债：房贷 411 万、车贷 15.6 万、装修贷 15.4 万、信用卡分期 6.6 万，合计 \textbf{448.6 万}。净值 \textbf{负 8.4 万}。

\textbf{孙桂芳}，41 岁，在蚌埠一家医院做护士，月薪 1 万 4，到手 1 万 1。

资产：没有贷款的房子 63 万、存款 14.8 万、指数基金 6.2 万、公积金 4.7 万，合计 \textbf{88.7 万}。负债 \textbf{0}。净值 \textbf{88.7 万}。

饭桌上所有人都会说李振宇有钱。

真正分胜负的是断粮那天。李振宇每月刚性支出是房贷 2 万 1 千 8、车贷 4300、装修贷 3800、信用卡 2200、基本生活 6000，加起来 \textbf{3 万 8 千 1}。账上 3.2 万，撑 \textbf{25 天}。孙桂芳每月刚性支出 3400，可动用的金融资产 21 万，撑 \textbf{61 个月}。

\end{egbox}

收入是流量，净值是存量。流量决定你现在过得怎么样，存量决定你扛得住什么。失业、生病、行业塌方，考的全是存量。

\section{一家奶茶店的三个月}

个人有一张表就够。公司要三张。

\begin{tablewrap}
\begin{xltabular}{\linewidth}{>{\raggedright\arraybackslash}p{0.20\linewidth}XX}
\toprule
\thead{报表} & \thead{它是什么} & \thead{回答的问题} \\
\midrule
\textbf{资产负债表} & 某一瞬间的\textbf{照片} & 这一刻我有什么、欠什么、剩什么 \\
\textbf{利润表} & 一段时间的\textbf{录像} & 这三个月我到底赚了还是赔了 \\
\textbf{现金流量表} & 真金白银的\textbf{进出账} & 这三个月钱包里的钱是多了还是少了 \\
\bottomrule
\end{xltabular}
\end{tablewrap}

2026 年 3 月，王雪梅在政务区一个写字楼底商开了家奶茶店，招牌两个字是她自己在 iPad 上描的。自己拿 22.6 万，向她妈借 18 万，说好三年后还，年息 3\%。

借条是她妈手写的，压在家里的玻璃板底下。老太太很讲原则。

这 40.6 万花掉：装修 14.8 万、设备 9.6 万、押一付三的房租 5.4 万、首批原料 3.2 万，剩 7.6 万压箱底。

\begin{bsheet}
  \bsside{开业当天 · 资产}{%
    \bsrow{装修（待摊）}{14.8 万}
    \bsrow{设备}{9.6 万}
    \bsrow{预付租金}{5.4 万}
    \bsrow{存货}{3.2 万}
    \bsrow{现金}{7.6 万}
    \bsrowtotal{合计}{40.6 万}
  }
  \bsside{负债 + 净值}{%
    \bsrow{欠她妈}{18 万}
    \bsrow{净值（王雪梅的本钱）}{22.6 万}
    \bsrowtotal{合计}{40.6 万}
  }
\end{bsheet}

4 到 6 月生意不错。卖了 2.87 万杯，均价 14.5 元，营业额 41.6 万。其中 4.8 万是给两家公司做团建供应的挂账，说好 9 月结账。这在会计上叫\term{应收账款}{accounts receivable}\index{076@应收账款}，是"赚了但还没收到"的钱。实际进账现金 36.8 万。

\begin{egbox}{举个栗子}
\textbf{利润表（4 到 6 月，万元）}：营业收入 41.6，扣掉原料 12.4、人工 9.3、房租 5.4、水电杂项 2.35、折旧摊销 1.71、利息 0.135，剩下 \textbf{利润 10.31}。

\textbf{现金流量表（4 到 6 月，万元）}：收到现金 36.8，付出原料 10.8（还欠供应商 1.6）、人工 9.3、水电 2.35、利息 0.135，房租这季度实付 0（开业前就付掉了），\textbf{现金净增 14.22}。期末现金 7.6 加 14.22，等于 \textbf{21.82}。

同一家店，同样三个月：利润 10.31 万，现金多了 14.22 万。差的 3.91 万从哪来？

房租 5.4，上季度就付了现金，这季度才计成本。折旧 1.71，根本没花钱。欠供应商的 1.6，还没付。三项加起来再减去挂账没收到的 4.8，正好 \textbf{3.91}。

分毫不差。

\end{egbox}

第二幕在 7 月。要续交下季度房租 5.4 万，还供应商 1.6 万，暑期旺季还得压 5.8 万的货。账上 21.82 万，撑得住。

可要是她这时候接了个 28.5 万的挂账大单，为此提前备 11.4 万的料，利润表会更好看，而现金在往外流。再接两单这样的"好生意"，她就得去借经营贷发工资了。

利润是一种观点，现金是一个事实。

利润取决于你怎么确认收入、怎么摊销、折旧年限估几年，里面装着一堆判断。账上还剩多少钱，没得商量。

\hlmark{很多公司死的时候，利润表还是正的。}

折旧那 1.71 万里，藏着整套会计最反直觉的一个概念。

王雪梅那台 9 万 6 的设备能用 5 年。她不能在买它的那个月把 9 万 6 全记成成本，那会让当月巨亏、后面四年多虚增利润，两头都不真实。正确做法是每年"消耗"1 万 9 千 2。这笔钱没有一分从账户流出去，可它确实是成本：设备真的在变旧，五年后她真的还得再掏一次钱。

\term{折旧}{depreciation}\index{084@折旧}就是把一次性的现金流出，摊到它真正被消耗的那些年里。

关键在于，折旧年限是一个判断，不是一个事实。

\begin{egbox}{举个栗子}
一家公司买了 137 亿的服务器。按 3 年折旧，每年成本 \textbf{45.7 亿}；按 6 年折旧，每年 \textbf{22.8 亿}。

同样的机器、同样的电费、同样的收入，账面利润差 \textbf{22.9 亿}，现金流一分钱没变。

这不是学术问题。2025 年全球数据中心资本开支同比增长约 57\%，规模越大，"用几年"这个假设的杠杆越大。\hlwarm{如果服务器实际 3 年就被新一代芯片挤下去，却按 6 年摊，那些省下来的成本迟早要一次性认回来。}AI 服务器该按几年折旧，是 2026 年最热的一场会计争论，目前没有定论。第 24 章会回到这里。

\end{egbox}

\section{三个口袋}

\begin{deepbox}{进阶：可以跳过}
把复式记账的逻辑放大到整个国家，会得到一个让很多人不舒服的结论。

把一个经济体切成三个口袋：政府、私人部门（居民加企业）、对外部门。因为\textbf{每一笔支出都是另一个人的收入}，三个口袋的收支之和必然为零。整理一下就是：

\textbf{政府赤字 = 私人部门盈余 + 对外部门赤字}。写成符号是 G-T =（S-I）+（M-X）。

\hlmark{这是会计恒等式，不是政策主张。}它只陈述一个算术事实：三个口袋不可能同时净盈余。政府压缩赤字，在其他条件不变时，要么私人部门积累金融资产的速度变慢，要么贸易顺差得变大，二者必居其一。

所以"政府必须像家庭一样量入为出"这句话，在会计上是错的。一个家庭省钱会变富，可\textbf{所有人同时省钱会让所有人的收入一起下降}，因为你的支出就是我的收入。经济学管这叫节俭悖论。

\textbf{这绝不意味着政府可以无限借债。}真正的约束不是"会破产"。一个只用本币举债、又拥有自己央行的政府，名义违约是政策选择而非技术必然。真正的约束有四条：\\ ① \textbf{通胀}。政府花钱争夺的是钢筋、工人和电力这些真实资源，资源不够时，多出来的钱只会变成物价。\\ ② \textbf{外部平衡}。赤字带动进口和资本外流，汇率承压。\\ ③ \textbf{币种和期限结构}。用外币借的债，你的央行印不出来；短期债要不断滚动，市场一旦不买单，利率瞬间跳升。历史上真正的主权债务危机几乎全出在这一条。\\ ④ \textbf{利息负担}。付息会挤占其他所有支出。

2026 年 8 月，美国国债突破 40 万亿美元，10 年期国债收益率升至 2025 年 1 月以来最高。市场在给期限风险和通胀风险重新定价，不是在给破产风险定价。围绕这套逻辑（MMT 是其中最激进的版本）学界吵得很凶，本书不站队，只请你分清哪一部分是会计，哪一部分是政策判断。

\end{deepbox}

\section{两副该摘的眼镜}

\begin{mythbox}{常见误解}
\mvfrow{误解}{"我有一套两百多万的房子，所以我有两百多万。"}
\mvfrow{事实}{你有的是\textbf{房产估值减去房贷余额}。而这两个数字的性质完全不同：\textbf{估值有弹性，负债没有弹性。}}
王雪梅家挂牌 232 万，房贷余额 176.4 万，账面净值 55.6 万。真卖一次看看：谈到 214 万成交，中介费按 1.5\% 算 3.21 万，各项税费 2.4 万，实收 208.39 万，还完贷款剩 \textbf{31.99 万}。

账面 55.6 万，落袋 32 万，缩水 \textbf{42.4\%}。这还没算从挂牌到过户的几个月。急着用钱的时候，时间本身就是最贵的成本。这叫\term{变现摩擦}{liquidity friction}\index{006@变现摩擦}。

再看负债那栏：176.4 万是精确到分的、白纸黑字的、每月必须付的。它不会因为行情不好就体谅你。\textbf{估值可以商量，还款日不能商量。}几乎所有个人财务危机都长在这个不对称上。这不是说买房错了，是说别把账面净值当成能动用的钱。第 19 章会教你做一张诚实的表。

\end{mythbox}

\begin{mythbox}{常见误解}
\mvfrow{误解}{"国家欠这么多债，跟一个家庭欠了一屁股债一样，早晚要还不起。"}
\mvfrow{事实}{这个类比在四个地方都不成立。但它的反面同样错：政府债务确实有硬约束，只是约束不叫"破产"。}
\textbf{寿命不同。}家庭必须在退休前还清，国家没有生命周期，只要增长率不长期低于利率，债务可以一直滚下去。\\ \textbf{支出的性质不同。}你花的钱不是你自己的收入，政府花的钱恰恰是私人部门的收入。\\ \textbf{定价权不同。}你没法影响自己的贷款利率，拥有独立央行的政府对本国利率有相当大的影响力。\\ \textbf{债主不同。}你的债主是外人，一国国债的相当一部分是欠本国居民、本国养老金和本国银行的，那既算负债，也算国民的资产。

\textbf{以上四条都不构成"可以随便借"的理由。}天花板是通胀、汇率和币种结构。历史上出事的国家，多数不是因为债务占 GDP 的比例太高，是因为借了自己印不出来的货币。

\end{mythbox}

\bookbreak

回到主线。钱是记账，而这一章给的是读账的语法。往后无论谈银行、央行、债券、股票、杠杆还是危机，都只是在看几张资产负债表怎么互相咬合、怎么一起膨胀、又怎么在某一天同时收缩。

王雪梅那张表上，90.3 万的净值里有 55.6 万锁在一套挂了四个月没卖动的房子里。

\begin{takeaway}{本章一句话}
资产 = 负债 + 净值。这个恒等式管你、管你楼下的奶茶店、管你的银行、也管你的国家；在这套语法里，一个人的金融资产必然是另一个人的负债。

\begin{itemize}
  \item 净值才是财富，收入只是流量。提前还贷不改变净值，涨薪存下来才改变净值。
  \item 利润是一种观点，现金是一个事实。折旧、应收账款和预付租金会让两者长期背离。
  \item 部门恒等式是会计不是政策。政府债务的真正约束是通胀、外部平衡和币种结构。
  \item 估值有弹性，负债没有弹性。所有财务危机都长在这个不对称上。
\end{itemize}

\end{takeaway}

\begin{bridgebox}
\textbf{下一章}有了这副眼镜，可以问那个最根本的问题了：把所有花哨名词剥干净之后，金融这门生意到底在交易什么？
\end{bridgebox}

